\section{MARCO TEÓRICO}

\subsection{Antecedentes de la investigación}
\parrafoII{Aguirre, J., Bonifaz, J. L., Chang, V., Flores, A., Lucich, I., Vega, A. y Mallma, D. (2025), en su estudio titulado Evaluaciones de impacto de asociaciones público-privadas en el Perú: Un trabajo conjunto para la generación de evidencia en 6 sectores de infraestructura y servicios públicos [1], se plantearon como objetivo primordial evaluar el efecto de diversas intervenciones de infraestructura en indicadores de eficiencia y bienestar social. La investigación se desarrolló bajo un enfoque cuantitativo de tipo aplicado, empleando una metodología basada en diseños experimentales y cuasiexperimentales, tales como el de diferencia en diferencias y el emparejamiento por puntaje de propensión. La población abarcó múltiples sectores a nivel nacional, utilizando una muestra estratificada de beneficiarios y grupos de control, donde la recolección de datos se efectuó mediante bases de datos administrativas y encuestas de hogares procesadas con software estadístico avanzado.Los resultados principales demostraron que la implementación de estándares técnicos rigurosos y el monitoreo constante incrementaron significativamente la eficiencia operativa y redujeron los márgenes de error en la prestación de servicios. Los hallazgos revelaron que la complejidad de los entornos de ejecución influyó directamente en el rendimiento de los proyectos, encontrando una correlación entre la calidad de la infraestructura y la optimización de recursos. El estudio concluyó que la evidencia basada en datos es fundamental para validar la estabilidad de cualquier sistema implementado. Este antecedente resultó altamente relevante para la presente investigación, ya que ofreció un marco metodológico sólido sobre el uso de diseños experimentales para medir el impacto de una variable independiente sobre el rendimiento. Asimismo, la vinculación entre la complejidad de los escenarios y la precisión de los resultados en el estudio de Aguirre et al. permitió fortalecer la hipótesis de que variables externas afectan directamente la eficiencia computacional y la fidelidad de los algoritmos desarrollados.}

\parrafoII{El Observatorio Nacional de Prospectiva (2024), en su artículo técnico titulado Potenciación e implementación de la IA en la generación de nuevas tecnologías [2], analizó la importancia estratégica de integrar la inteligencia artificial para impulsar el desarrollo tecnológico. El estudio se planteó como objetivo principal examinar las capacidades de la IA para procesar grandes volúmenes de datos, realizar análisis predictivos y automatizar tareas complejas en diversos sectores industriales. La investigación se desarrolló bajo un enfoque cualitativo-prospectivo de tipo aplicado y documental, empleando una metodología de análisis de tendencias y revisión de marcos regulatorios internacionales. La población del estudio abarcó las estrategias nacionales de IA de países líderes (como Canadá, Singapur y el Reino Unido), utilizando como técnica de recolección el análisis de contenido de reportes gubernamentales y guías de innovación actualizadas al año 2024.Los resultados principales destacaron que la implementación de la IA actúa como un motor de eficiencia industrial y crecimiento económico, permitiendo la creación de nuevas soluciones tecnológicas que transforman industrias enteras. Los hallazgos subrayaron que la sinergia entre la IA y la generación de tecnologías optimiza la gestión de recursos y mejora la capacidad de respuesta ante problemas técnicos complejos. El estudio concluyó que para una implementación exitosa se requieren marcos éticos y una inversión constante en investigación que valide la estabilidad de las nuevas herramientas. Este antecedente es relevante para la investigación actual, ya que fundamenta la importancia de la IA en la generación de componentes lógicos y algoritmos eficientes. Al vincularse con el problema general, este documento respalda la hipótesis de que la IA no solo automatiza la creación de código, sino que su eficacia y rendimiento dependen de la infraestructura estratégica y la complejidad de los proyectos tecnológicos donde se aplique.}

\parrafoII{Aparicio-Gómez, O. Y. y Aparicio-Gómez, W. O. (2024), en su artículo de reflexión titulado Innovación educativa con sistemas de aprendizaje adaptativo impulsados por Inteligencia Artificial [3], analizaron cómo los algoritmos de aprendizaje automático personalizan procesos complejos mediante el procesamiento de datos a gran escala. El estudio tuvo como objetivo examinar la integración de sistemas adaptativos basados en IA para optimizar la eficacia y los resultados en entornos de aprendizaje. La investigación se desarrolló bajo un enfoque cualitativo de tipo reflexivo y documental, empleando una metodología de análisis crítico de literatura científica contemporánea. La población consistió en una selección de marcos teóricos y aplicaciones tecnológicas de machine learning y deep learning, utilizando como técnica de recolección el mapeo sistemático de fuentes académicas indexadas y reportes de innovación tecnológica actualizados al año 2024. Los resultados principales destacaron que los algoritmos de IA, al ser entrenados con grandes volúmenes de datos, logran generar itinerarios precisos y respuestas en tiempo real que superan los modelos estáticos tradicionales. Los hallazgos subrayaron que la arquitectura de estos sistemas permite una retroalimentación constante, lo cual es determinante para la estabilidad del desempeño en tareas de alta complejidad. El estudio concluyó que la infraestructura tecnológica y la calidad de los datos son pilares fundamentales para que la IA mantenga niveles óptimos de motivación y eficacia. Este antecedente es relevante para la investigación actual, ya que proporciona un sustento teórico sobre la capacidad de los algoritmos de IA para adaptarse a escenarios variables. Al vincularse con el problema de la precisión y el rendimiento del código, este estudio respalda la hipótesis de que la calidad del "entrenamiento" y la estructura del algoritmo influyen directamente en la fidelidad de los resultados generados ante la complejidad del escenario físico planteado.}

\parrafoII{Chaparro Aguilar, E., Arias Escobar, J. y Quispe Capajaña, M. (2025), en su artículo original titulado Uso de inteligencia artificial generativa en contextos académicos: percepción de estudiantes universitarios [4], se propusieron como objetivo analizar la percepción y la frecuencia de uso de herramientas de IA generativa para el desarrollo de actividades académicas. La investigación se realizó bajo un enfoque cuantitativo, de tipo descriptivo y transversal, empleando una metodología basada en el método deductivo. La población estuvo conformada por estudiantes de ingeniería de universidades públicas del sur del Perú, de la cual se extrajo una muestra de 216 participantes mediante un muestreo no probabilístico por conveniencia, utilizando como técnica de recolección la encuesta a través de un cuestionario estructurado y validado por juicio de expertos. Los resultados principales revelaron que la mayoría de los estudiantes utiliza la IA para la generación y revisión de contenidos técnicos, destacando una percepción positiva sobre la utilidad de estas herramientas para mejorar la productividad. Los hallazgos indicaron que, si bien existe un uso frecuente, persisten desafíos relacionados con la verificación de la precisión de los resultados entregados por los modelos. El estudio concluyó que la integración de la IA en la formación de ingenieros es irreversible y requiere el desarrollo de competencias críticas para validar la veracidad de la información generada. Este antecedente es relevante para la investigación actual, ya que contextualiza el perfil de los usuarios (estudiantes de ingeniería) que interactúan con la IA para producir código. Al vincularse con el problema de la precisión física y el rendimiento, este estudio aporta evidencia sobre la confianza depositada en estas herramientas, reforzando la necesidad de evaluar experimentalmente si dicha confianza se traduce en resultados técnicos fieles y eficientes en escenarios de alta complejidad.}

\parrafoII{El Instituto Nacional de Defensa de la Competencia y de la Protección de la Propiedad Intelectual - INDECOPI (2025), mediante la Resolución N° 000062-2025-GEG/INDECOPI y su documento normativo adjunto titulado Lineamientos para el Uso Ético de la Inteligencia Artificial en el Instituto Nacional de Defensa de la Competencia y de la Protección de la Propiedad Intelectual - Indecopi [5], estableció un marco regulatorio para la implementación de sistemas basados en IA. El estudio se planteó como objetivo primordial definir principios y directrices técnicas y éticas que garanticen el uso responsable, transparente y seguro de la inteligencia artificial dentro de las funciones públicas. La investigación se desarrolló bajo un enfoque normativo y técnico de tipo aplicado, empleando una metodología de análisis legal y prospectivo basada en la Ley N° 31814. La población abarcó a todos los servidores civiles y proveedores de la entidad, utilizando como técnica de recolección la revisión de estándares internacionales de gobernanza de datos y auditoría de algoritmos. Los resultados principales consistieron en la estructuración de un Inventario de IA y la determinación de niveles de riesgo para cada sistema desarrollado o implementado. Los hallazgos subrayaron que la supervisión humana y la explicabilidad del código son requisitos indispensables para asegurar la estabilidad y veracidad de los procesos automatizados. El estudio concluyó que cualquier desarrollo tecnológico basado en IA debe ser sometido a una evaluación de impacto que valide su integridad y el cumplimiento de estándares de seguridad de la información. Este antecedente es fundamental para la investigación actual, ya que proporciona el marco regulatorio peruano vigente que exige la validación y el control de calidad sobre el código generado por IA. Su relevancia se vincula directamente con los objetivos de medir el rendimiento y la precisión, pues establece que la eficiencia no puede estar disociada de la supervisión técnica y la reducción de riesgos algorítmicos en escenarios de alta complejidad.}

\parrafoII{Villasana Lopez, M. (2025), en su tesis titulada Inteligencia artificial generativa en la integridad académica en los estudiantes de ingeniería de sistemas de una Universidad de Huancavelica en 2024 [6], se planteó como objetivo determinar la relación entre el uso de la inteligencia artificial generativa y la integridad académica. La investigación se desarrolló bajo un enfoque cuantitativo, de tipo aplicado y nivel correlacional, empleando un diseño no experimental y transversal. La población estuvo constituida por estudiantes de la facultad de ingeniería de sistemas, de la cual se seleccionó una muestra probabilística de 136 participantes, utilizando como técnica de recolección la encuesta y el análisis de logs de procesos mediante un script especializado para rastrear el uso de herramientas como ChatGPT. Los resultados principales demostraron una correlación significativa entre la frecuencia de interacción con modelos de lenguaje y la modificación de los procesos de producción de trabajos técnicos. Los hallazgos revelaron que el uso de la IA se concentra en la resolución de problemas algorítmicos y la generación de código fuente, donde la percepción de eficacia del estudiante influye en la profundidad del análisis crítico realizado. El estudio concluyó que la integración de estas tecnologías en ingeniería de sistemas demanda mecanismos de validación técnica que aseguren la autoría y la precisión del software resultante. Este antecedente es relevante para la investigación actual, ya que aborda directamente a la población de ingeniería de sistemas y el uso de herramientas generativas para tareas de codificación. Su vinculación con el problema de la precisión y el rendimiento radica en que proporciona evidencia sobre cómo el comportamiento del usuario y la confianza en la herramienta pueden condicionar la calidad técnica final del producto desarrollado.}

\parrafoII{Gorosito, A. R., Carbonell, A. E., Marcipar, L. y Gieco, L. (2025), en su artículo de investigación titulado Inteligencia artificial en la formación de ingenieros [7], evaluaron la incorporación de la inteligencia artificial en el proceso de enseñanza y aprendizaje de estudiantes de ingeniería electromecánica. El estudio se planteó como objetivo primordial identificar los errores, ventajas y limitaciones de la IA como colaborador inteligente en la resolución de problemas técnicos complejos. La investigación se desarrolló bajo un enfoque cualitativo de tipo descriptivo y exploratorio, empleando una metodología basada en entrevistas estructuradas y observación de campo durante el periodo 2024-2025. La población estuvo conformada por docentes y estudiantes de cátedras integradoras, utilizando como técnicas de recolección de datos la entrevista directa y el análisis de la interacción grupal con herramientas de IA en laboratorios de computación. Los resultados principales demostraron que, si bien la IA actúa como un facilitador del aprendizaje autónomo, presenta deficiencias críticas en el análisis contextual de problemas de ingeniería específicos. Los hallazgos revelaron que las soluciones generadas por los modelos a menudo carecen de precisión cuando no son supervisadas con sentido común y rigor técnico, detectando errores en la adaptación de respuestas a situaciones físicas reales. El estudio concluyó que la tecnología no sustituye la capacidad analítica, sino que requiere una validación constante para asegurar la estabilidad de las soluciones propuestas. Este antecedente es altamente relevante para la investigación actual, ya que aborda la problemática de la precisión y los errores técnicos en el código generado. Al vincularse con los objetivos de este estudio, el trabajo de Gorosito et al. refuerza la hipótesis de que la complejidad del escenario físico pone a prueba las limitaciones de la IA, exigiendo una evaluación experimental de su rendimiento y fidelidad algorítmica.}

\parrafoII{Elnashar, A., Moundas, M., Schmidt, D. C., Spencer-Smith, J. y White, J. (2024), en su artículo de investigación titulado Evaluating the Performance of LLM-Generated Code for ChatGPT-4 and AutoGen Along with Top-Rated Human Solutions [8], se propusieron como objetivo realizar un análisis comparativo de la eficacia y el rendimiento del código generado por modelos de lenguaje de gran escala frente a soluciones humanas de alta calificación. La investigación se desarrolló bajo un enfoque cuantitativo de tipo aplicado, empleando una metodología de experimentación controlada y benchmarking. La población consistió en una selección de problemas algorítmicos complejos extraídos de Stack Overflow, utilizando como muestra 100 soluciones generadas por ChatGPT-4 y 50 casos de prueba dinámicos para el framework AutoGen, donde la recolección de datos se efectuó mediante entornos de prueba dinámicos que midieron el éxito de ejecución y métricas de rendimiento en tiempo de ejecución. Los resultados principales demostraron que las soluciones generadas por IA alcanzan tasas de éxito superiores al 90 por ciento, resultando competitivas e incluso superiores a las mejores soluciones humanas en términos de corrección lógica. Los hallazgos revelaron que, a pesar de la alta eficacia, los modelos presentan dificultades específicas al manejar tipos de datos de punto flotante y casos de borde en algoritmos de ordenamiento, lo que afecta la robustez de la solución. El estudio concluyó que la ingeniería de prompts y el uso de múltiples agentes colaborativos optimizan la calidad del software, aunque persiste la necesidad de atender errores de precisión en tipos de datos complejos. Este antecedente es fundamental para la investigación actual, ya que proporciona un referente directo sobre la evaluación del rendimiento computacional de la IA generativa. Su relevancia se vincula estrechamente con el problema de la precisión física y el rendimiento, al aportar evidencia empírica sobre cómo la complejidad del problema algorítmico influye en la aparición de errores técnicos y en la eficiencia del código resultante.}

\parrafoII{Azbarka, S., Sabbagh, M., Mughrabi, A. y Roth, I. (2024), en su estudio de investigación titulado Evaluation of the Code Quality Generated by Generative AI [9], se plantearon como objetivo evaluar y diferenciar la efectividad de las principales herramientas de inteligencia artificial generativa mediante el examen de métricas clave de calidad del código. La investigación se desarrolló bajo un enfoque cuantitativo y comparativo de tipo aplicado, empleando una metodología estructurada en dos etapas de benchmarking. La población consistió en los modelos de lenguaje más avanzados del mercado, seleccionando una muestra compuesta por ChatGPT, GitHub Copilot, Claude y Gemini, los cuales fueron sometidos a la resolución de 40 problemas de programación de diversa dificultad extraídos de la plataforma LeetCode, utilizando Python como lenguaje de implementación y entornos de ejecución controlados para la recolección de datos. Los resultados principales demostraron que ChatGPT obtuvo la puntuación más alta en la mayoría de las métricas de calidad, logrando una tasa de corrección superior en comparación con sus competidores. Los hallazgos revelaron que, aunque herramientas como Copilot y Claude mostraron un rendimiento sólido, existen variaciones significativas en la eficiencia del tiempo de ejecución y la legibilidad del código dependiendo de la complejidad del problema planteado. El estudio concluyó que la elección de la herramienta de IA influye directamente en la calidad técnica del software producido, sugiriendo que los desarrolladores deben seleccionar el modelo basándose en la naturaleza específica de la tarea algorítmica. Este antecedente es de vital importancia para la investigación actual, ya que proporciona un marco comparativo directo entre las herramientas de IA (variable independiente) y su impacto en la precisión y el rendimiento (variables dependientes). Su vinculación con el problema general fortalece la necesidad de evaluar experimentalmente cómo estos mismos modelos se comportan al enfrentarse a la complejidad de simulaciones físicas, más allá de los problemas de lógica puramente algorítmica.}

\parrafoII{Cao, Y., Lai, S., Huang, J., Zhang, Y., Lawrence, Z., Bhakta, R., Cao, M., Tsai, C., Zhou, Z., Zhao, Y., Liu, H., Marinoni, A., Arefiev, A., Yu, R. y Thomas, I. F. (2026), en su artículo de investigación titulado SimulCost: A Cost-Aware Benchmark for Automating Physics Simulations with LLMs [10], se propusieron como objetivo introducir y evaluar un nuevo marco de referencia para medir la precisión y el costo computacional de los modelos de lenguaje de gran escala (LLM) al automatizar simulaciones físicas complejas. La investigación se desarrolló bajo un enfoque cuantitativo de tipo aplicado y tecnológico, empleando una metodología de benchmarking a gran escala que integró 12 simuladores de diversas áreas como la dinámica de fluidos y la mecánica de sólidos. La población del estudio consistió en los modelos de lenguaje más avanzados (incluyendo GPT-5, Claude y Llama), utilizando una muestra de 2,916 tareas de ronda única y 1,900 de ajuste iterativo, donde la recolección de datos se realizó mediante la ejecución de simulaciones y la medición analítica del tiempo de cómputo y el error de precisión. Los resultados principales demostraron que los LLM de última generación alcanzan tasas de éxito de entre el 46\% y 64\% en la generación inicial, pero este rendimiento cae drásticamente por debajo del 54\% cuando se exigen niveles altos de precisión física. Los hallazgos revelaron que, aunque el modo de ajuste iterativo mejora la exactitud, los modelos de IA resultan ser entre 1.5 y 2.5 veces más lentos que los métodos de optimización tradicionales, evidenciando una ineficiencia en el manejo de recursos computacionales. El estudio concluyó que las conjeturas iniciales de la IA son poco fiables para tareas científicas que requieren alta fidelidad, subrayando la necesidad de equilibrar el costo de simulación con la precisión obtenida. Este antecedente es crítico para la investigación actual, ya que aborda directamente la relación entre la complejidad del escenario físico y el rendimiento computacional. Al vincularse con los objetivos de este estudio, el trabajo de Cao et al. proporciona evidencia empírica de que la IA presenta degradaciones en la precisión y un uso ineficiente de recursos al enfrentarse a simulaciones de alta complejidad, validando la relevancia de evaluar estas variables experimentalmente.}

\parrafoII{Wang, Y., Wang, R., Wang, Y., Liang, X., Koto, F., Baldwin, T., Liang, X. y Li, H. (2026), en su estudio de investigación titulado SimuScene: Training and Benchmarking Code Generation to Simulate Physical Scenarios [11], se plantearon como objetivo evaluar la capacidad de los modelos de lenguaje de gran escala (LLM) para representar y simular escenarios físicos precisos a través de la generación de código. La investigación se desarrolló bajo un enfoque cuantitativo de tipo aplicado y tecnológico, empleando una metodología de entrenamiento y evaluación sistemática que abarcó cinco dominios de la física y 52 conceptos físicos. La población consistió en una selección de modelos contemporáneos de IA, utilizando una muestra de 7,659 escenarios físicos generados automáticamente, de los cuales 334 fueron verificados por humanos como conjunto de prueba, empleando técnicas de recolección de datos basadas en pipelines automáticos y recompensas visuales mediante modelos de lenguaje-visión (VLM). Los resultados principales demostraron que incluso los modelos más avanzados alcanzan una tasa de éxito de apenas el 21.5\% en la simulación precisa de escenarios físicos, lo que evidencia la alta dificultad de la tarea. Los hallazgos revelaron errores recurrentes como la penetración de objetos en límites sólidos y la inconsistencia entre la descripción textual y la dinámica física renderizada. El estudio concluyó que el uso de aprendizaje por refuerzo con retroalimentación visual mejora significativamente la generación de código orientado a la física, aunque la brecha de precisión sigue siendo amplia. Este antecedente es de suma relevancia para la investigación actual, ya que aborda directamente la problemática de la fidelidad física en el código generado por IA. Al vincularse con los objetivos de medir la precisión y el rendimiento, este trabajo valida la hipótesis de que los modelos actuales presentan limitaciones críticas al enfrentar la complejidad de leyes físicas reales, subrayando la importancia de realizar pruebas experimentales exhaustivas en entornos de simulación.}
\subsection{Bases teóricas}